\documentclass[11pt,a4paper]{article}
\usepackage[margin=21mm,top=22mm,bottom=20mm]{geometry}
\usepackage{amsmath,amssymb}
\usepackage{booktabs}
\usepackage{tabularx}
\usepackage{array}
\usepackage{xcolor}
\usepackage[most]{tcolorbox}
\usepackage{titlesec}
\usepackage{fancyhdr}
\usepackage{enumitem}
\usepackage[colorlinks=true,linkcolor=black,urlcolor=black]{hyperref}

\definecolor{mid}{HTML}{5B6B72}
\definecolor{pale}{HTML}{EEF3F4}
\definecolor{acc}{HTML}{0D5D70}
\setlength{\parskip}{4pt}
\setlength{\parindent}{0pt}
\titleformat{\section}{\large\bfseries\color{acc}}{\thesection}{0.6em}{}
\titleformat{\subsection}{\normalsize\bfseries}{\thesubsection}{0.5em}{}
\setlist{itemsep=1pt,topsep=2pt}
\newcolumntype{L}{>{\raggedright\arraybackslash}X}
\newtcolorbox{eqbox}{colback=white,colframe=acc,boxrule=0.5pt,arc=2pt,left=6pt,right=6pt,top=3pt,bottom=3pt}
\newtcolorbox{result}{colback=pale,colframe=pale,boxrule=0pt,arc=2pt,left=8pt,right=8pt,top=4pt,bottom=4pt}
\newcommand{\dfus}{\Delta_{\mathrm{fus}}}
\newcommand{\aoh}{a_{\mathrm{OH}}}
\newcommand{\file}[1]{\texttt{\small #1}}
\newcommand{\aice}{\ln a_w^{\mathrm{ice}}}

\pagestyle{fancy}\fancyhf{}
\fancyhead[L]{\footnotesize\color{mid}Cryo-SLE calculations workbook}
\fancyhead[R]{\footnotesize\color{mid}S. Chowdhary, 25B0347 --- CL207}
\fancyfoot[C]{\footnotesize\thepage}
\renewcommand{\headrulewidth}{0.3pt}

\begin{document}

{\LARGE\bfseries\color{acc} Calculations workbook}\par\vspace{2pt}
{\large Every thermodynamic equation in the project, worked by hand on one example}\par\vspace{4pt}
{\small\color{mid} Siddhant Chowdhary (25B0347), CL207 optional project. Companion to the report and the handbook.}\par
\vspace{6pt}

\begin{result}
\textbf{How to read this.} Each section gives one equation, the numbers going in, every intermediate step, and the result,
which is then checked against data. All numbers here were produced by \file{analysis/worked\_calculations.py}. That script
does each calculation twice, once line by line from the formula as written here and once with the project's own code,
and stops with an error if the two disagree. So this workbook cannot drift from the code that made the report's results.
Rerun it with \file{python analysis/worked\_calculations.py}.
Numbers are rounded for printing; the calculations carry full precision.
\end{result}

\section*{Map: equation $\rightarrow$ where it is used $\rightarrow$ worked example}
\small
\begin{tabularx}{\textwidth}{@{}l L L r@{}}
\toprule
& Equation & Used for & Example result\\
\midrule
1 & Ice line: $\aice(T)$ & Every freezing calculation & $a_w=<<ice_aw>>$ at $-15\,^\circ$C\\
2 & Composition: wt\% $\rightarrow$ $x$, $m$ & Every calculation & $x_w=<<g_xw>>$ (20 wt\% glycerol)\\
3 & Ideal SLE: $\ln x_w=\aice$ & Baseline & $<<id_TC>>\,^\circ$C (data $<<g_crc>>$)\\
4 & Margules $\ln\gamma_w$, fit of $\aoh$ & The one fitted constant & $\aoh=<<AOH>>$ J/mol\\
5 & Non-ideal SLE & Fitted solute & $<<no_TC>>\,^\circ$C (data $<<g_crc>>$)\\
6 & Same, $A=n_\mathrm{OH}\aoh$ & Blind tests & EG 28 wt\%: $<<e_ours>>$ (data $<<e_crc>>$)\\
7 & NaCl: $\ln a_w=-2m M_w\phi(m)$ & Salt in every mixture & saline: $<<s_T>>\,^\circ$C\\
8 & Additive mixing in $\ln a_w$ & Ternary tests, recipes & gly 25 + NaCl 5: $<<t_T>>\,^\circ$C\\
9 & Freeze concentration (water balance) & Salt in unfrozen liquid & saline, $-15\,^\circ$C: $<<fc_wt>>$ wt\% NaCl\\
10 & Same, with glycerol & Recipe criterion 1 & 6.5 wt\% gly: $<<c_wt>>$ wt\% NaCl\\
11 & Cell: $a_w^{\mathrm{in}}=a_w^{\mathrm{ice}}$ & Salt inside the cell & sorbitol $<<in_wta>>$ / gly $<<in_wtb>>$ wt\%\\
12 & Osmotic coefficient at 25 $^\circ$C & Sorbitol test at 25 $^\circ$C & $\phi$: model $<<b_phim>>$, data $<<b_phi>>$\\
13 & Gibbs--Helmholtz $a=h-Ts$ & Post-hoc explanation & $h=<<gh0_h>>$ J/mol per OH\\
14 & Sparing ratio $-\mathrm{d}w_g/\mathrm{d}w_s$ & Design result & $<<r_out_fd>>$ / $<<r_in_fd>>$ / $<<r_unf_fd>>$\\
\bottomrule
\end{tabularx}
\normalsize

\section*{Constants and the one fitted number}
\begin{tabularx}{\textwidth}{@{}l l L@{}}
\toprule
Symbol & Value & Meaning / source\\
\midrule
$R$ & $<<R>>$ J mol$^{-1}$ K$^{-1}$ & gas constant\\
$T_m$ & $<<TM>>$ K & melting point of ice at 1 atm\\
$\dfus H$ & $<<DH>>$ J mol$^{-1}$ & enthalpy of fusion of ice at $T_m$\\
$\Delta C_p$ & $<<DCP>>$ J mol$^{-1}$ K$^{-1}$ & $C_p(\text{liquid})-C_p(\text{ice})$ near 0 $^\circ$C\\
$M_w$ & $<<Mw>>$ g mol$^{-1}$ & molar mass of water\\
$\aoh$ & $<<AOH>>$ J mol$^{-1}$ per OH & \textbf{the only parameter fitted to polyol data} (glycerol, CRC, 2--40 wt\%)\\
$a,b,c$ & $<<na>>,\ <<nb>>,\ <<nc>>$ & NaCl osmotic coefficient, fitted to the NaCl--water ice line only\\
\bottomrule
\end{tabularx}

Useful derived number, the cryoscopic constant: $K_f=RT_m^2M_w/\dfus H=<<R>>\times<<TM>>^2\times0.018015/<<DH>>=<<Kf>>$ K kg mol$^{-1}$
(textbook value 1.86, so the constants are consistent).

%-----------------------------------------------------------------
\section{The ice line: how much water activity ice allows at $T$}
\begin{eqbox}
\[
\aice(T)=\underbrace{-\frac{\dfus H}{R}\Big(\frac1T-\frac1{T_m}\Big)}_{\text{term 1: latent heat}}
\;\underbrace{-\;\frac{\Delta C_p}{R}\Big(\frac{T_m}{T}-1+\ln\frac{T}{T_m}\Big)}_{\text{term 2: heat-capacity correction}}
\tag{1}\]
\end{eqbox}
\textbf{Where it comes from.} At equilibrium $\mu_w^{\text{ice}}=\mu_w^{\text{liq}}=\mu_w^{\text{L},0}+RT\ln a_w$, so
$RT\ln a_w=-\dfus G(T)$. $\dfus G$ is known only at $T_m$ (where it is zero), so it is carried to $T$ along
warm ice $\rightarrow$ melt at $T_m$ $\rightarrow$ cool liquid (Sandler ch.~11; handbook \S2.2--2.3).

\textbf{Example: $T=-15\,^\circ$C $=<<iceT>>$ K.}
\begin{align*}
\tfrac1T-\tfrac1{T_m} &= \tfrac1{<<iceT>>}-\tfrac1{<<TM>>}=<<ice_inv>>\ \text{K}^{-1}\\
\text{term 1} &= -\,<<ice_DHR>>\ \text{K}\times<<ice_inv>>=<<ice_t1>>\\
\tfrac{T_m}{T}-1+\ln\tfrac{T}{T_m} &= <<ice_r>>\\
\text{term 2} &= -\,<<ice_DCPR>>\times<<ice_r>>=<<ice_t2>>\\
\aice &= <<ice_t1>><<ice_t2>>=\mathbf{<<ice_ln>>},\qquad a_w^{\text{ice}}=e^{<<ice_ln>>}=\mathbf{<<ice_aw>>}
\end{align*}
\begin{result}
Any liquid touching ice at $-15\,^\circ$C has $a_w=<<ice_aw>>$, \emph{whatever is dissolved in it}. The $\Delta C_p$ term is
$<<ice_pct2>>$\% of the total here: small but not negligible at cryo temperatures, so it is kept everywhere.
\end{result}

%-----------------------------------------------------------------
\section{Composition: weight percent to mole fraction and molality}
\textbf{Example: 20 wt\% glycerol} ($M=<<Mg>>$ g/mol), basis 100 g of solution:
\begin{align*}
n_g&=20/<<Mg>>=<<g_ng>>\ \text{mol},\qquad n_w=80/<<Mw>>=<<g_nw>>\ \text{mol}\\
x_w&=\frac{n_w}{n_w+n_g}=\mathbf{<<g_xw>>},\quad x_g=<<g_xs>>,\qquad m_g=\frac{n_g}{0.080\ \text{kg}}=<<g_m>>\ \text{mol/kg}
\end{align*}
Polyols use mole fractions (they enter the excess Gibbs energy model); NaCl uses molality (its $\phi(m)$ is written per kg water).

%-----------------------------------------------------------------
\section{Ideal freezing point (Raoult, $\gamma_w=1$)}
\begin{eqbox}
\[ \ln x_w=\aice(T)\qquad\text{solve for } T. \tag{2}\]
\end{eqbox}
Eq.~(2) cannot be solved for $T$ in closed form (because of term 2), so it is solved numerically. Define
$g(T)=\ln x_w-\aice(T)$ and use Newton's method, $T_{k+1}=T_k-g(T_k)/g'(T_k)$, starting 1 K below $T_m$:

\begin{center}\small
\begin{tabular}{c c c}
\toprule
$k$ & $T_k$ (K) & $g(T_k)$\\
\midrule
0 & <<id_T0>> & <<id_g0>>\\
1 & <<id_T1>> & <<id_g1>>\\
2 & <<id_T2>> & <<id_g2>>\\
3 & <<id_T3>> & <<id_g3>>\\
\bottomrule
\end{tabular}
\end{center}
(The project code uses Brent's method, which is safer for hard cases; it gives the same root.)
\begin{result}
Ideal: $T_f=<<id_T>>$ K $=\mathbf{<<id_TC>>\,^\circ}$\textbf{C}. CRC data: $<<g_crc>>\,^\circ$C. The ideal model is
$0.7$ K too warm, so real glycerol depresses the freezing point \emph{more} than Raoult's law: $\gamma_w<1$.
Dilute shortcut $\Delta T=-K_f m=-<<Kf>>\times<<g_m>>=<<id_dil>>$ K, already off at this concentration.
\end{result}

%-----------------------------------------------------------------
\section{The activity model and fitting $\aoh$}
\begin{eqbox}
\begin{gather*}
\frac{G^{ex}}{RT}=x_w\sum_i\alpha_i x_i,\qquad \alpha_i=\frac{A_i}{RT},\qquad A_i=n_{\mathrm{OH},i}\,\aoh\\
\Rightarrow\quad \ln\gamma_w=(1-x_w)\sum_i\alpha_i x_i,\qquad \ln\gamma_i=\alpha_i x_w-x_w\sum_j\alpha_jx_j \tag{3}
\end{gather*}
\end{eqbox}
For one solute, eq.~(3) is the one-constant Margules model, $RT\ln\gamma_w=A\,x_s^2$. Both $\ln\gamma$'s come from
differentiating $nG^{ex}$, so they satisfy Gibbs--Duhem automatically. $A<0$ means water--solute attraction (hydrogen bonds), so $\gamma_w<1$.

\textbf{What one data point says.} At the CRC freezing point of 20 wt\% glycerol, $T=<<f_T>>$ K:
\begin{align*}
\ln a_w &= \aice(<<f_T>>)=<<f_lnaw>> &&\text{(eq.~1: ice fixes $a_w$)}\\
\ln\gamma_w &= \ln a_w-\ln x_w=<<f_lnaw>>-(<<g_lnxw>>)=<<f_lngw>>,\quad \gamma_w=<<f_gw>>\\
\aoh &= \frac{RT\ln\gamma_w}{n_\mathrm{OH}\,x_g^2}=\frac{<<R>>\times<<f_T>>\times(<<f_lngw>>)}{3\times<<g_xs>>^2}=<<f_a>>\ \text{J/mol per OH}
\end{align*}
\textbf{The fit.} Doing this properly means choosing $\aoh$ to minimise $\sum(T_f^{\text{model}}-T_f^{\text{data}})^2$ over all
$<<fit_n>>$ CRC glycerol points (2--40 wt\%) at once (\file{analysis/fit\_literature.py}):
\begin{result}
$\aoh=\mathbf{<<AOH>>}$ \textbf{J mol$^{-1}$ per OH}; RMSE $<<fit_rmse>>$ K, worst point $<<fit_max>>$ K. The single-point
estimate ($<<f_a>>$) is within 2\% of it. This number is then frozen: it is never refitted on any test data.
\end{result}

%-----------------------------------------------------------------
\section{Non-ideal freezing point of the fitted solute}
\begin{eqbox}
\[ \ln x_w+\frac{A}{RT}(1-x_w)^2=\aice(T),\qquad A=3\aoh=<<no_A>>\ \text{J/mol for glycerol}\tag{4}\]
\end{eqbox}
Newton on $g(T)=\ln x_w+\frac{A}{RT}x_g^2-\aice(T)$ with $x_w=<<g_xw>>$:
\begin{center}\small
\begin{tabular}{c c c}
\toprule
$k$ & $T_k$ (K) & $g(T_k)$\\
\midrule
0 & <<no_T0>> & <<no_g0>>\\
1 & <<no_T1>> & <<no_g1>>\\
2 & <<no_T2>> & <<no_g2>>\\
3 & <<no_T3>> & <<no_g3>>\\
\bottomrule
\end{tabular}
\end{center}
At the root: $\alpha=A/RT=<<no_alpha>>$; $\ln\gamma_w=\alpha x_g^2=<<no_lng>>$ ($\gamma_w=<<no_gw>>$);
glycerol's own $\ln\gamma_g=\alpha x_w^2=<<no_lngs>>$, a large negative deviation.
\begin{result}
$T_f=<<no_T>>$ K $=\mathbf{<<no_TC>>\,^\circ}$\textbf{C} vs data $<<g_crc>>\,^\circ$C. A correction of under 1\% in $\gamma_w$
moves $T_f$ by 0.7 K, which is why the activity coefficient is the whole problem.
\end{result}

%-----------------------------------------------------------------
\section{A blind prediction: ethylene glycol, 28 wt\%}
Same eq.~(4), nothing refitted: ethylene glycol has 2 OH, so $A=2\aoh=<<e_A>>$ J/mol. Composition ($M=<<e_Me>>$):
$x_w=<<e_xw>>$.
\begin{center}\small
\begin{tabular}{l c c}
\toprule
Model & $T_f$ ($^\circ$C) & error vs CRC (K)\\
\midrule
Ideal, eq.~(2) & <<e_id>> & <<e_err_id>>\\
Per-OH, eq.~(4) & \textbf{<<e_ours>>} & \textbf{<<e_err>>}\\
CRC data & <<e_crc>> & ---\\
\bottomrule
\end{tabular}
\end{center}
Across all 2--40 wt\% EG points this gives the Table~1 result of the report (RMSE 0.83 K vs ideal 2.49, UNIFAC 3.99).
The same calculation with $n_\mathrm{OH}=6$ is what overstates sorbitol (section~12).

%-----------------------------------------------------------------
\section{Salt: NaCl through an osmotic coefficient}
\begin{eqbox}
\[ \ln a_w^{\mathrm{NaCl}}=-\nu\, m\, M_w\,\phi(m),\quad \nu=2,\qquad
\phi(m)=1-\frac{a\sqrt m}{1+\sqrt m}+b\,m+c\,m^2 \tag{5}\]
\end{eqbox}
The $\sqrt m$ term has the Debye--H\"uckel shape (ion clouds make $\phi<1$ when dilute); $b,c$ capture the rise at high $m$.
$a,b,c$ were fitted to the NaCl--water ice line only (17 CRC points).

\textbf{Example: isotonic saline, 0.9 wt\%} ($M=<<s_Mn>>$):
\begin{align*}
m&=\frac{0.9/<<s_Mn>>}{0.0991\ \text{kg}}=<<s_m>>\ \text{mol/kg},\quad \sqrt m=<<s_sqrt>>\\
\phi&=1-\frac{<<na>>\times<<s_sqrt>>}{1+<<s_sqrt>>}+<<nb>>\times<<s_m>>+<<nc>>\times<<s_m>>^2=<<s_phi>>\\
\ln a_w&=-2\times<<s_m>>\times0.018015\times<<s_phi>>=<<s_lnaw>>
\end{align*}
\begin{result}
Setting $\ln a_w=\aice(T)$ gives $T_f=\mathbf{<<s_T>>\,^\circ}$\textbf{C}. Ideal-ion shortcut $-K_f\nu m=<<s_dil>>$ K;
with $\phi$, $<<s_dilphi>>$ K. Saline freezes at about $-0.54\,^\circ$C, the standard value.
\end{result}

%-----------------------------------------------------------------
\section{Mixtures: additive in $\ln a_w$}
\begin{eqbox}
\[ \ln a_w^{\mathrm{mix}}=\underbrace{\ln x_w+(1-x_w)\textstyle\sum_i\alpha_ix_i}_{\text{polyols, eq. 3 (water + polyols only)}}
\;+\;\underbrace{\left(-2mM_w\phi(m)\right)}_{\text{NaCl, eq. 5}} \;=\;\aice(T)\tag{6}\]
\end{eqbox}
Pre-registered on 21 Oct, before any ternary data: no salt--polyol cross term. Bower \& Robinson's (1963) ternary data put the
neglected term at 0.6--5\% of the osmolality.

\textbf{Example: the committed lab sample, glycerol 25 + NaCl 5 wt\%} (70 g water):
$x_w$(water + glycerol) $=<<t_xw>>$; $m_{\text{NaCl}}=<<t_m>>$ mol/kg, $\phi=<<t_phi>>$. At the root:
\[
\ln a_w=\underbrace{<<t_poly>>}_{\text{glycerol}}\;+\;\underbrace{(<<t_salt>>)}_{\text{NaCl}}=<<t_ln>>=\aice(T_f)
\]
\begin{result}
$T_f=\mathbf{<<t_T>>\,^\circ}$\textbf{C}, identical to the value committed in \file{predictions/lab\_set\_v3\_glycerol\_nacl.csv}
($<<t_comm>>$) before any measurement. Glycerol supplies two-thirds of the water-activity drop.
\end{result}

%-----------------------------------------------------------------
\section{Freeze concentration: what is left liquid at $-15\,^\circ$C}
Ice is pure water, so as water freezes out, every solute stays in a shrinking liquid. Freezing stops when the liquid's
$a_w$ has fallen to the ice value. The unknown is the mass of liquid water left, $W$; the moles of solute are fixed.
\begin{eqbox}
\[ \ln a_w^{\mathrm{mix}}\big(n_1,n_2,\dots;\,W\big)=\aice(T)\quad\text{solve for } W \tag{7}\]
\end{eqbox}
\textbf{Example 1: saline only.} With only NaCl, eq.~(7) is $-2mM_w\phi(m)=\aice(<<iceT>>)=<<fc_target>>$, an equation in $m$ alone:
\begin{align*}
m^{*}&=<<fc_m>>\ \text{mol/kg}\ (\phi=<<fc_phi>>)\\
\text{unfrozen water}&=W/W_0=m_0/m^{*}=<<s_m>>/<<fc_m>>=\mathbf{<<fc_u>>\,\%}\\
\text{NaCl in the liquid}&=\frac{m^{*}M}{1+m^{*}M}=\mathbf{<<fc_wt>>\ wt\%}
\end{align*}
96\% of the water is ice; the cells sit in near-brine. ($m^{*}$ is below the top of the fitted NaCl range, $<<fc_mmax>>$ mol/kg.)

%-----------------------------------------------------------------
\section{Freeze concentration with glycerol (recipe criterion 1)}
\textbf{Example 2: glycerol 6.5 wt\% in the saline.} Basis 100 g: $n_g=<<c_ng>>$ mol, $n_{\text{NaCl}}=<<c_nn>>$ mol,
$W_0=<<c_W0>>$ g. Unknown $W$ (g of liquid water). For a trial $W$:
$x_w=\frac{W/M_w}{W/M_w+n_g}$ and $m=\frac{n_{\text{NaCl}}}{W/1000}$; put both into eq.~(6) and solve eq.~(7) for $W$ (Brent's method):
\begin{align*}
W&=<<c_W>>\ \text{g}\quad(\text{unfrozen water } <<c_u>>\%)\\ x_w&=<<c_xw>>,\quad m_{\text{NaCl}}=<<c_mn>>,\quad m_g=<<c_mg>>\ \text{mol/kg}\\
\ln a_w&=\underbrace{<<c_poly>>}_{\text{glycerol}}+\underbrace{(<<c_salt>>)}_{\text{NaCl}}=<<fc_target>>=\aice\quad\checkmark\\
\text{NaCl in liquid}&=\frac{0.9}{<<c_W>>+6.5+0.9}=\mathbf{<<c_wt>>\ wt\%}\quad(\text{glycerol } <<c_gwt>>\ \text{wt\%})
\end{align*}
\begin{result}
The ice still sets the same $a_w$; glycerol just takes over most of the job salt was doing. 6.5 wt\% is the least
glycerol that holds NaCl at the 4.41 wt\% damage limit (Lovelock 1953) at $-15\,^\circ$C, by this criterion.
\end{result}

%-----------------------------------------------------------------
\section{Inside the cell}
Water crosses the membrane freely, so at (slow-cooling) equilibrium
\begin{eqbox}
\[ \ln a_w^{\mathrm{in}}\Big(\tfrac{n_c}{W_c},\ m_g\Big)=\ln a_w^{\mathrm{out}}=\aice(T)\tag{8}\]
\end{eqbox}
Assumptions: the cell's own solutes $n_c$ are represented as NaCl at the isotonic level ($m_0=<<in_ms0>>$ mol/kg);
glycerol penetrates (same molality inside and out); sorbitol and NaCl stay out. Unknown: cell water $W_c/W_{c,0}$.

\textbf{(a) Sorbitol 9 wt\% outside.} Nothing permeant, so the inside holds only $n_c$, and eq.~(8) is \emph{the same equation
as the saline case}:
\[ m^{\mathrm{in}}=m^{*}=<<fc_m>>,\qquad W_c/W_{c,0}=<<in_ms0>>/<<fc_m>>=<<in_Wa>>,\qquad \text{salt inside}=\mathbf{<<in_wta>>\ wt\%}. \]
Outside, the sorbitol has cut NaCl to $<<in_out_sorb>>$ wt\%, but the inside is as salty as with no protectant at all.

\textbf{(b) Glycerol 6.5 wt\%.} Glycerol enters at $m_g=<<c_mg>>$ mol/kg (from section 10), sharing the water-activity drop:
\[ W_c/W_{c,0}=<<in_Wb>>,\qquad m^{\mathrm{in}}=<<in_ms0>>/<<in_Wb>>=<<in_msb>>,\qquad \text{salt inside}=\mathbf{<<in_wtb>>\ wt\%}. \]
\begin{result}
Because ice fixes $a_w$ on \emph{both} sides of the membrane, only a solute that gets in can lower the salt inside. This is
Pegg's ``must penetrate'' (2007), derived from $\mu$-equality rather than assumed. It holds for any impermeant solute,
not only in this simplified cell.
\end{result}

%-----------------------------------------------------------------
\section{Sorbitol at 25 $^\circ$C: osmotic coefficient}
Away from ice, isopiestic data report the osmotic coefficient, defined by
\begin{eqbox}
\[ \ln a_w=-\phi\, m\, M_w\quad(\nu=1)\qquad\Leftrightarrow\qquad \phi=-\frac{\ln a_w}{mM_w}\tag{9}\]
\end{eqbox}
\textbf{Example: Bower \& Robinson (1963), $m=<<b_m>>$ mol/kg, measured $\phi=<<b_phi>>$.} $x_s=<<b_xs>>$.
\begin{align*}
\text{model: }\ln a_w&=\ln x_w+\frac{6\aoh}{R\times298.15}x_s^2=<<b_lnaw>>\ \Rightarrow\ \phi=\mathbf{<<b_phim>>}\\
\text{ideal: }\phi&=-\ln x_w/(mM_w)=<<b_phiid>>
\end{align*}
Turn it round: the data's $\ln a_w=-<<b_phi>>\times<<b_m>>\times0.018015=<<b_lnawdat>>$ implies
$a=\frac{(\ln a_w-\ln x_w)RT}{6x_s^2}=\mathbf{<<b_a>>}$ J/mol per OH, about a quarter of glycerol's $<<AOH>>$.
\begin{result}
The per-OH model predicts $\phi-1$ about 6$\times$ too large here; ideal is closer. The constant does not transfer to sorbitol.
Section 13 separates how much of this is temperature and how much is the molecule.
\end{result}

%-----------------------------------------------------------------
\section{Gibbs--Helmholtz: splitting $a$ into $h$ and $s$ (post hoc)}
\begin{eqbox}
\[ a(T)=h-Ts\qquad\Rightarrow\qquad s=-\frac{a(T_2)-a(T_1)}{T_2-T_1},\qquad h=a(T_1)+T_1s\tag{10}\]
\end{eqbox}
$a$ is an interaction \emph{free} energy, so it has an enthalpy part (hydrogen-bond strength) and an entropy part (ordering of water).
Sucrose is the one solute with good data at two temperatures (CRC freezing points; Robinson \& Stokes 1949 at 25 $^\circ$C).
Implied $a$ at each point is found as in section 4, and medians are compared in matched molality windows:

\begin{center}\small
\begin{tabular}{l c c c c c}
\toprule
Window (mol/kg) & $a(T_1)$, $T_1$ & $a(298.15)$ & $s=-\Delta a/\Delta T$ & $h=a_1+T_1s$ & $T_1 s$\\
\midrule
<<gh0_win>> & $<<gh0_a1>>$, <<gh0_T1>> K & $<<gh0_a2>>$ & $-\frac{<<gh0_a2>>-(<<gh0_a1>>)}{<<gh0_dT>>}=<<gh0_s>>$ & $<<gh0_h>>$ & $<<gh0_Ts>>$\\
<<gh1_win>> & $<<gh1_a1>>$, <<gh1_T1>> K & $<<gh1_a2>>$ & $-\frac{<<gh1_a2>>-(<<gh1_a1>>)}{<<gh1_dT>>}=<<gh1_s>>$ & $<<gh1_h>>$ & $<<gh1_Ts>>$\\
\bottomrule
\end{tabular}
\end{center}
Units: $a,h$ in J/mol per OH; $s$ in J/mol/K per OH. (The script's unrounded values are $h=<<gh0_hcode>>$ and $<<gh1_hcode>>$.)
\begin{result}
$h\approx-6$ to $-8$ kJ/mol per OH (the size of a water hydrogen bond), with an entropy penalty of about $-15$ to $-20$
J/mol/K. Warming makes $-Ts$ bigger, so the attraction weakens by about 20\% from the freezing line to 25 $^\circ$C. That
is far too little to explain sorbitol ($-1800$ for sucrose vs $<<b_a>>$ for sorbitol at the same 25 $^\circ$C), so sorbitol's
failure is a molecule effect.
\end{result}

%-----------------------------------------------------------------
\section{The design number: glycerol saved per gram of sorbitol}
For each sorbitol dose $w_s$, find the least glycerol $w_g$ that meets a criterion at $-15\,^\circ$C (sections 10--11 inside
a root-finder). The sparing ratio is the slope $-\mathrm{d}w_g/\mathrm{d}w_s$, read here over the first 2 wt\% of sorbitol:

\begin{center}\small
\begin{tabular}{l c c c c}
\toprule
Criterion & $w_g$ at $w_s=0$ & $w_g$ at $w_s=2$ & ratio $(w_g^0-w_g^2)/2$ & linear fit (range)\\
\midrule
NaCl outside $\le$ 4.41 wt\% & <<r_out_0>> & <<r_out_2>> & \textbf{<<r_out_fd>>} & <<r_out>> (<<r_out_rng>>)\\
salt inside $\le$ 4.41 wt\% & <<r_in_0>> & <<r_in_2>> & \textbf{<<r_in_fd>>} & <<r_in>> (<<r_in_rng>>)\\
unfrozen water $\ge$ 25\% & <<r_unf_0>> & <<r_unf_2>> & \textbf{<<r_unf_fd>>} & <<r_unf>> (<<r_unf_rng>>)\\
\bottomrule
\end{tabular}
\end{center}
\begin{result}
Positive means sorbitol saves glycerol; negative means it costs glycerol. With all three criteria at once, the binding constraint switches
from ``unfrozen water'' to ``salt inside'' and the best mix is sorbitol $<<r_opt_s>>$ + glycerol $<<r_opt_g>>$ wt\%, against
$<<r_alone>>$ wt\% glycerol alone ($<<r_saving>>$\% less glycerol). Whether sorbitol helps depends on which injury mechanism
limits survival, and the model puts a number on each case.
\end{result}

\section*{Reproduce}
\small
\begin{tabular}{@{}l l@{}}
All values in this workbook & \file{python analysis/worked\_calculations.py} (asserts hand = code)\\
Fit of $\aoh$ & \file{analysis/fit\_literature.py}\\
Freeze concentration, cell, criteria & \file{analysis/freeze\_concentration\_cells.py}\\
Sorbitol at 25 $^\circ$C & \file{analysis/sorbitol\_25C.py}\\
Gibbs--Helmholtz & \file{analysis/posthoc\_per\_molecule.py}\\
Everything & \file{python run\_all.py}\\
\end{tabular}

\end{document}
